% 设计总览：多播-聚合流与组播 credit（fig:design-overview）
% (a) 一份 normal credit 的组播闭环；(b) 一个 channel 的多棵 subchannel 树
\begin{minipage}[b]{0.56\linewidth}
\centering
\begin{tikzpicture}[x=1cm,y=1cm,
  box/.style={draw,rounded corners=1.5pt,minimum height=0.5cm,
    inner sep=3pt,font=\scriptsize},
  acc/.style={box,fill=ArborLightGray!60},
  ep/.style={box,fill=ArborBlue!12},
  credit/.style={-stealth,ArborBlue,semithick,dashed},
  up/.style={-stealth,ArborOrange,semithick},
  lab/.style={font=\tiny,inner sep=1pt}]
  \node[ep]  (R)  at ( 0.0,4.55) {responder};
  \node[acc] (SP) at ( 0.0,3.15) {加速器};
  \node[acc] (L1) at (-1.9,1.75) {加速器};
  \node[acc] (L2) at ( 1.9,1.75) {加速器};
  \node[ep]  (q1) at (-2.7,0.35) {$q_1$};
  \node[ep]  (q2) at (-1.1,0.35) {$q_2$};
  \node[ep]  (q3) at ( 1.9,0.35) {$q_3$};
  \node[lab,anchor=west] at ( 0.62,3.38) {slot 7};
  \node[lab,anchor=east] at (-2.52,1.98) {slot 3};
  \node[lab,anchor=west] at ( 2.52,1.98) {slot 5};
  % 下行：credit 组播副本，逐级推进 allocator 并压栈
  \draw[credit] (R)  to[bend right=18]
    node[lab,left=2pt]{credit $(k,s_1)$} (SP);
  \draw[credit] (SP) to[bend right=18] (L1);
  \draw[credit] (SP) to[bend right=18] (L2);
  \draw[credit] (L1) to[bend right=18]
    node[lab,left=2pt]{栈 $[7,3]$} (q1);
  \draw[credit] (L1) to[bend right=18] (q2);
  \draw[credit] (L2) to[bend right=18]
    node[lab,right=2pt]{栈 $[7,5]$} (q3);
  % 上行：request 按栈逐级聚合并累计贡献数
  \draw[up] (q1) to[bend right=18] (L1);
  \draw[up] (q2) to[bend right=18] (L1);
  \draw[up] (L1) to[bend right=18] node[lab,right=2pt]{$\times2$} (SP);
  \draw[up] (q3) to[bend right=18] (L2);
  \draw[up] (L2) to[bend right=18] node[lab,left=2pt]{$\times1$} (SP);
  \draw[up] (SP) to[bend right=18]
    node[lab,right=2pt]{$\times3$，完成样本} (R);
  % 图例
  \draw[credit] (-3.30,-0.55) -- (-2.85,-0.55);
  \node[lab,anchor=west] at (-2.80,-0.55) {credit 下行：分配 slot、压位置栈};
  \draw[up] ( 0.55,-0.55) -- ( 1.00,-0.55);
  \node[lab,anchor=west] at ( 1.05,-0.55) {request 上行：按栈聚合};
\end{tikzpicture}
\par\smallskip{\small (a) 一份 normal credit 的组播闭环}
\end{minipage}\hfill
\begin{minipage}[b]{0.40\linewidth}
\centering
\begin{tikzpicture}[x=1cm,y=1cm,
  box/.style={draw,rounded corners=1.5pt,minimum height=0.5cm,
    inner sep=3pt,font=\scriptsize},
  ep/.style={box,fill=ArborBlue!12},
  sw/.style={box,fill=ArborLightGray!60},
  t1/.style={ArborBlue,semithick},
  t2/.style={ArborOrange,semithick,dashed},
  lab/.style={font=\tiny,inner sep=1pt}]
  \node[sw] (SA) at (-0.9,2.05) {$S_A$};
  \node[sw] (SB) at ( 0.9,2.05) {$S_B$};
  \node[ep] (Rb) at (-2.4,0.40) {R};
  \node[ep] (b1) at (-0.8,0.40) {$q_1$};
  \node[ep] (b2) at ( 0.8,0.40) {$q_2$};
  \node[ep] (b3) at ( 2.4,0.40) {$q_3$};
  \foreach \h in {Rb,b1,b2,b3}{
    \draw[t2] (SB) -- (\h);
  }
  \foreach \h in {Rb,b1,b2,b3}{
    \draw[t1] (SA) -- (\h);
  }
  \node[lab,anchor=east,ArborBlue]   at (-1.38,2.42) {树 $s_1$};
  \node[lab,anchor=west,ArborOrange] at ( 1.38,2.42) {树 $s_2$};
  \node[lab,align=center] at (0,3.00)
    {offset：$k\to s_1$，\ $k{+}1\to s_2$，\ $k{+}2\to s_1$，$\cdots$};
  \node[lab,align=center] at (0,-0.30)
    {每棵树独立维护 $r_s$/$W_s$ gate；gate 关闭的树不取得新 offset};
\end{tikzpicture}
\par\smallskip{\small (b) 一个 channel 的多棵 subchannel 树}
\end{minipage}
